% Source printed pp. 327-343 / PDF pp. 335-351.
\section{Gerbe, 刚性化与 Picard 概形}\label{src:7.4}
\begin{epigraph}
大多数超椭圆曲线都没有点.
\epigraphauthor{Don Zagier}
\end{epigraph}
Gerbe 在叠论中起核心作用. 本书需要知道: (a) 剩余 gerbe 确实是 gerbe (命题~\bookref{7.4.39}), 从而这一术语名副其实; (b) 从 Picard 叠到 Picard 概形的映射 $\underline{\mathcal P\!ic}_X\to\underline{\Pic}_X$, 以及 (c) 从稳定向量丛叠到其粗模空间的映射 $\mathcal B\!un_{r,d}^s(C)\to M_{r,d}^s(C)$, 都是带 $\mathbb G_m$-带的 gerbe (推论~\bookref{9.2.26}).
\subsection{挠子}\label{src:7.4.1-subsection}
对群层 $G$, $G$-挠子是局部同构于 $G$ 的集合层. 当 $G$ 可由 fppf 群概形表示时, 这一概念等价于主 $G$-丛.
\begin{definition}[挠子]\label{src:7.4.1}
设 $\mathcal S$ 是景, $G$ 是 $\mathcal S$ 上的 (未必 Abel 的) 群层. $\mathcal S$ 上的\emph{$G$-挠子} 是 $\mathcal S$ 上的集合层 $P$, 配备 $G$ 的左作用 $\sigma:G\times P\to P$, 满足
\begin{enumerate}[(1)]
\item 对每个对象 $T\in\mathcal S$, 存在覆盖 $\{T_i\to T\}$, 使每个 $i$ 都有 $P(T_i)\neq\varnothing$.
\item 作用映射 $(\sigma,p_2):G\times P\to P\times P$ 是同构.
\end{enumerate}
若 $T\in\mathcal S$ 是对象, $G$ 是限制景 $\mathcal S/T$ (例~\bookref{3.2.12}) 上的群层, 则 $T$ 上的 $G$-挠子按定义是 $\mathcal S/T$ 上的 $G$-挠子.
\end{definition}
$G$-挠子的态射是层的 $G$-等变态射. 若 $P$ 与 $G$ 是 $G$-等变同构的, 则称 $G$-挠子 $P$ \emph{平凡}.
\begin{exercise}\label{src:7.4.2}
证明 $G$-挠子的任意态射都是同构.
\end{exercise}
\begin{example}\label{src:7.4.3}
设 $\mathcal X$ 是景 $\mathcal S$ 上的叠, $a,b\in\mathcal X$ 是 $S\in\mathcal S$ 上的对象. 同构层 $\underline{\operatorname{Isom}}_S(a,b)$ 在由预复合给出的作用下是 $\underline{\operatorname{Aut}}(a)$ 的挠子.
\end{example}
给定态射 $f:T'\to T$ 及 $T$ 上的 $G$-挠子 $P$, 限制 $P|_{T'}$ 是 $\mathcal S/T'$ 上的层, 在 $T'$-概形 $S$ 上的截面为 $P(S)$; 它自然是 $T'$ 上的 $G$-挠子.
\begin{exercise}\label{src:7.4.4}
设 $\mathcal S$ 是具有终对象 $S$ 的景, $G$ 是 $\mathcal S$ 上的群层.
\begin{enumerate}[(a)]
\item 证明公理 (1) 等价于 $P\to S$ 是层的满态射.
\item 若 $P$ 是 $G$-挠子, 证明 $S$ 同构于商层 $P/G$.
\item 证明 $G$-挠子 $P$ 平凡, 当且仅当结构态射 $P\to S$ 存在截面 $s:S\to P$.
\item 证明 $\mathcal S$ 上带 $G$ 左作用的集合层 $P$ 是 $G$-挠子, 当且仅当存在覆盖 $\{S_i\to S\}$ 及 $G|_{S_i}$-挠子的同构 $P|_{S_i}\cong G|_{S_i}$.
\end{enumerate}
\end{exercise}
\begin{example}[主 \texorpdfstring{$G$}{G}-丛]\label{src:7.4.5}
对 fppf 群概形 $G\to S$, fppf 拓扑中的 $G$-挠子与主 $G$-丛 (定义~\bookref{7.3.7}) 的范畴等价. 为说明这一点, 注意若 $P\to T$ 是 $S$-概形 $T$ 上的主 $G$-丛, 则由于代数空间是 fppf 拓扑中的层 (推论~\bookref{7.3.3}), $G\times_S T$ 与 $P$ 都是 $(\Sch/T)_{\mathrm{fppf}}$ 上的层. $P\to T$ 被 fppf 态射 $P\to T$ 平凡化, 由习题~\bookref{7.4.4}(d), 它是 $G$-挠子. 反之, 给定 $(\Sch/T)_{\mathrm{fppf}}$ 上的 $G$-挠子 $P$ 及使 $P$ 平凡化的 fppf 覆盖 $T'\to T$, 即 $P\times_T T'\cong G\times_T T'$, 则 $P\times_T T'$ 是代数空间, $P\times_T T'\to T'$ 是 fppf 可表示态射. 代数性的 fppf 判据 (\bookref{7.3.1}) 表明 $P$ 是代数空间, 故 $P\to T$ 是主 $G$-丛. 若 $G\to S$ 还光滑, 则平展拓扑中的 $G$-挠子与主 $G$-丛的范畴等价. 这是因为光滑态射平展局部有截面, 因而每个主 $G$-丛 $P\to T$ 都平展局部平凡.
\end{example}
\subsection{Gerbe}\label{src:7.4.2-subsection}
Gerbe 是挠子的 $2$-范畴推广. 挠子局部同构于群层 $G$, 而 gerbe 局部同构于分类叠 $BG$.
\begin{definition}[Gerbe]\label{src:7.4.6}
若景 $\mathcal S$ 上的叠 $\mathcal X$ 满足下列条件, 则称它为 \newterm{gerbe}{gerbe}:
\begin{enumerate}[(1)]
\item 对每个对象 $T\in\mathcal S$, 存在 $\mathcal S$ 中的覆盖 $\{T_i\to T\}$, 使每个纤维范畴 $\mathcal X(T_i)$ 非空.
\item 对 $T\in\mathcal S$ 上的对象 $x,y\in\mathcal X$, 存在覆盖 $\{T_i\to T\}$ 及每个 $i$ 上的同构 $x|_{T_i}\xrightarrow{\sim}y|_{T_i}$.
\end{enumerate}
若 $\mathcal X\to\mathcal S$ 有截面 $\mathcal S\to\mathcal X$, 则称 gerbe $\mathcal X$ \emph{平凡}. 当 $\mathcal S$ 有终对象 $S$ 时, 平凡性等价于 $\mathcal X(S)$ 非空.
\end{definition}
\begin{example}\label{src:7.4.7}
对景 $\mathcal S$ 上的群层 $G$, 推广定义~\bookref{7.3.8}, 将 $G$ 的分类叠定义为 $\mathcal S$ 上由对 $(P,T)$ 组成的叠 $BG$, 其中 $T\in\mathcal S$, $P$ 是 $\mathcal S/T$ 上的 $G$-挠子 (定义~\bookref{7.4.1}). 态射 $(P',T')\to(P,T)$ 是 $\mathcal S$ 中的态射 $T'\to T$ 及 $G$-挠子的同构 $P'\to P|_{T'}$ 的数据, 其中 $P|_{T'}$ 表示 $P$ 沿 $T'\to T$ 的限制. 分类叠 $BG$ 是 $\mathcal S$ 上的 gerbe, 因为 $T$ 上每个 $G$-挠子都局部同构于平凡 $G$-挠子 $G\times T$.
\end{example}
\begin{exercise}[Gerbe 局部是分类叠]\label{src:7.4.8}
设 $\mathcal S$ 是具有终对象 $S\in\mathcal S$ 的景, $\mathcal X$ 是 $\mathcal S$ 上的叠. 证明 $\mathcal X$ 是 gerbe, 当且仅当存在覆盖 $\{S_i\to S\}$ 及限制景 $\mathcal S/S_i$ (例~\bookref{3.2.12}) 上的群层 $G_i$, 使在 $\mathcal S/S_i$ 上有同构 $\mathcal X\times_S S_i\cong BG_i$.
\end{exercise}
\begin{exercise}\label{src:7.4.9}
设 $S$ 是概形, $\mathcal X$ 是 $(\Sch/S)_{\mathrm{fppf}}$ 上的 gerbe. 若对角 $\mathcal X\to\mathcal X\times_S\mathcal X$ 可表示, 证明 $\mathcal X$ 是代数叠.
\end{exercise}
用群层 $G$ 给 gerbe $\mathcal X$ 指定带, 是对每个对象 $x\in\mathcal X(T)$ 额外指定自然同构 $G(T)\to\operatorname{Aut}_T(x)$ 的数据.
\begin{definition}[带 \texorpdfstring{$G$}{G}-带的 gerbe]\label{src:7.4.10}
设 $G$ 是景 $\mathcal S$ 上的 Abel 层. \newterm{带 $G$-带的 gerbe}{banded $G$-gerbe} (或简称 $G$-gerbe) 是 $\mathcal S$ 上的 gerbe $\mathcal X$, 配备对每个对象 $x\in\mathcal X(T)$ 的层同构 $\psi_x:G|_T\to\underline{\operatorname{Aut}}_T(x)$. 要求对 $T$ 上每个同构 $\alpha:x\xrightarrow{\sim}y$, 图
\begin{equation}\label{src:7.4.11}
\begin{tikzcd}[ampersand replacement=\&]
\&G|_T\arrow[dl,"{\psi_x}"']\arrow[dr,"{\psi_y}"]\&\\
\underline{\operatorname{Aut}}_T(x)\arrow[rr,"{\operatorname{Inn}_\alpha}"']\&\&\underline{\operatorname{Aut}}_T(y)
\end{tikzcd}
\end{equation}
交换,\footnote{若 $G$ 是 (可能非 Abel 的) 群层, 则这一条件迫使 $G$ 为 Abel 群层.} 其中 $\operatorname{Inn}_\alpha(\tau)=\alpha\tau\alpha^{-1}$. 同构 $\psi_x$ 的数据称为 $\mathcal X$ 的 \newterm{$G$-带}{$G$-band}.\footnote{我们将在 \S\bookref{7.4.4-subsection} 中考虑更一般的带的概念.} 带 $G$-带的 gerbe 的态射是与 $G$-带相容的叠态射.
\end{definition}
\begin{remark}\label{src:7.4.12}
另一种理解 gerbe 的 $G$-带的方式是: 设 $\mathcal S/\mathcal X$ 是底层范畴为 $\mathcal X$ 的限制景, 对 $S\in\mathcal S$, $a\in\mathcal X(S)$ 的覆盖是 $S$ 的覆盖. 则惯性叠 $I_{\mathcal X}=\mathcal X\times_{\mathcal X\times\mathcal X}\mathcal X$ 是 $\mathcal S/\mathcal X$ 上的群层: 对 $a\in\mathcal X(S)$, 有 $I_{\mathcal X}(a)=\operatorname{Isom}_S(a)$. 相容条件 \bookeqref{7.4.11} 保证存在 $\mathcal S/\mathcal X$ 上的层同构 $\psi:G|_{\mathcal X}\to I_{\mathcal X}$.
\end{remark}
\begin{remark}\label{src:7.4.13}
带 $G$-带的 gerbe 可视为分类叠 $BG$ 下的挠子. 这是一个有用的理念性观点, 可以完全严格化, 参见 \cite[\S A.5]{CZ17}, 并有具体结果, 例如若 $\mathcal X\to X$ 是 $\mathbb G_m$-gerbe, 则 $D^b(\QCoh(\mathcal X))$ 带有典范的 $\mathbb Z$-分次.
\end{remark}
\begin{example}[平凡的带带 gerbe]\label{src:7.4.14}
若 $G$ 是具有终对象 $S$ 的景 $\mathcal S$ 上的 Abel 层, 则例~\bookref{7.4.7} 的分类叠 $BG$ 是带 $G$-带的 gerbe, 并且平凡. $\mathcal S$ 上带 $G$-带的 gerbe $\mathcal X$ 平凡, 当且仅当 $\mathcal X\cong BG$.
\end{example}
\begin{exercise}[Abel gerbe 所关联的带]\label{src:7.4.15}
设 $\mathcal S$ 是具有终对象 $S$ 的景, $\mathcal X$ 是 $\mathcal S$ 上的 Abel gerbe, 即对每个对象 $a\in\mathcal X(T)$, $\operatorname{Aut}_T(a)$ 都是 Abel 群. 证明存在 $\mathcal S$ 上的群层 $G$, 使 $\mathcal X$ 带 $G$-带.

\textit{提示:} 用 gerbe 的公理 (1) 找到覆盖 $\{X_i\to X\}$ 及元素 $a_i\in\mathcal X(X_i)$. 用公理 (2) 将层 $G_i:=\underline{\operatorname{Aut}}_{X_i}(a_i)$ 粘合为层 $G$.
\end{exercise}
\subsection{代数 gerbe}\label{src:7.4.3-subsection}
任意代数叠 $\mathcal Y$ 都有关联的大 fppf 景 $(\Sch/\mathcal Y)_{\mathrm{fppf}}$, 由 $\mathcal Y$ 上的概形组成: 底层范畴是 $\mathcal Y$, 对象 $y\in\mathcal Y(T)$ 的覆盖是 $T$ 的覆盖. 若 $\mathcal X\to\mathcal Y$ 是代数叠态射, 则由推论~\bookref{7.3.3}, $\mathcal X$ 是 $(\Sch/\mathcal Y)_{\mathrm{fppf}}$ 上的叠.
\begin{definition}[Gerbe 与带 \texorpdfstring{$G$}{G}-带的 gerbe]\label{src:7.4.16}
若 $\mathcal X$ 是大 fppf 景 $(\Sch/\mathcal Y)_{\mathrm{fppf}}$ 上的 gerbe, 则称代数叠态射 $\mathcal X\to\mathcal Y$ 为 \emph{gerbe}. 若 $G\to S$ 是交换 fppf 群概形, 则当 $S$ 上的代数叠态射 $\mathcal X\to\mathcal Y$ 是 $(\Sch/\mathcal Y)_{\mathrm{fppf}}$ 上的 gerbe, 且带群层 $G\times_S\mathcal Y$ 给出的带时, 称它为\emph{带 $G$-带的 gerbe} (或简称 $G$-gerbe). 若存在到代数空间的态射 $\mathcal X\to X$, 它是 gerbe (分别地, 带 $G$-带的 gerbe), 则称代数叠 $\mathcal X$ 为 gerbe (分别地, 带 $G$-带的 gerbe).
\end{definition}
代数空间 $X$ 上带 $G$-带的 gerbe $\mathcal X\to X$ 平凡, 当且仅当 $\mathcal X\cong BG\times_S X$.
\begin{proposition}\label{src:7.4.17}
设 $\mathcal X\to\mathcal Y$ 是代数叠态射.
\begin{enumerate}[(1)]
\item 态射 $\mathcal X\to\mathcal Y$ 是 gerbe, 当且仅当存在来自概形的 fppf 态射 $S\to\mathcal Y$ 及 fppf 群代数空间 $G\to S$, 使 $\mathcal X\times_{\mathcal Y}S\cong BG$.
\item 若 $\mathcal X\to\mathcal Y$ 是 gerbe, 则 $\mathcal X\to\mathcal Y$ 光滑.
\item 若 $G\to S$ 是 fppf 群概形, $\mathcal X\to X$ 是代数空间 $X$ 上带 $G$-带的 gerbe, 则存在平展覆盖 $X'\to X$, 使 $\mathcal X\times_X X'\cong X'\times_S BG$.
\end{enumerate}
\end{proposition}
\begin{proof}
对 (1) 的 ($\Rightarrow$), 由习题~\bookref{7.4.8}, 存在 fppf 态射 $S\to\mathcal Y$ 及 $S$ 上的群层 $G$, 使 $\mathcal X\times_{\mathcal Y}S\cong BG$. 由于 $\mathcal X\times_{\mathcal Y}S$ 是代数叠, 其对角可表示, 因而 $G$ 是代数空间. 反之, 若 $\mathcal X\times_{\mathcal Y}S\cong BG$, 则 gerbe 的公理 (1) 成立: 若 $a\in(\Sch/\mathcal Y)$ 是概形 $T$ 上的对象, 则 $S_T:=S\times_{\mathcal Y}T\to T$ 是 fppf 覆盖, 且 $\mathcal X\times_{\mathcal Y}S_T\cong BG_{S_T}$, 因而存在 $S_T$ 上的 $\mathcal X$ 的对象. 类似地, 对公理 (2), 若 $x_1,x_2\in\mathcal X$ 是 $y\in\mathcal Y(T)$ 上的对象, 则 $x_1,x_2$ 的拉回在 fppf 覆盖 $S_T\to T$ 下变得同构. (2) 由 (1) 得到, 因为 $BG\to S$ 光滑 (即使 $G\to S$ 不光滑!): 考虑复合 $S\to BG\to S$, 并使用光滑性在源上 fppf 局部 (命题~\bookref{3.1.27}), 即可得到这一点. 对 (3), 在光滑满射态射 $\mathcal X\to X$ 下作基变换后, $\mathcal X\to X$ 有截面. 选取光滑呈示 $U\to\mathcal X$, 则复合 $U\to\mathcal X\to X$ 光滑, 基变换 $\mathcal X\times_X U\to U$ 有截面. 结论由光滑态射平展局部有截面这一事实 (推论~\bookref{A.3.5}) 得到.
\end{proof}
\begin{exercise}\label{src:7.4.18}
证明代数叠态射 $\mathcal X\to\mathcal Y$ 是 gerbe, 当且仅当 $\mathcal X\to\mathcal Y$ 与 $\mathcal X\to\mathcal X\times_{\mathcal Y}\mathcal X$ 都是 fppf.
\end{exercise}
事实上, 具有平坦惯性的任意代数叠都是 gerbe. 下述结果将用于证明单向量丛的模叠是代数空间上的 $\mathbb G_m$-gerbe (命题~\bookref{9.2.23}).
\begin{proposition}\label{src:7.4.19}
设 $\mathcal X$ 是代数叠, $I_{\mathcal X}\to\mathcal X$ 是 fppf. 将给概形 $S$ 指派对象同构类集合 $\mathcal X(S)/\!\sim$ 的函子层化, 得到 $\Sch_{\mathrm{fppf}}$ 上的层 $X$. 则 $X$ 是代数空间, $\mathcal X\to X$ 是 gerbe.
\end{proposition}
\begin{proof}
为证明 $X$ 是代数空间, 只需证明 $\mathcal X\to X$ 是光滑可表示态射. 事实上, 光滑呈示 $U\to\mathcal X$ 诱导满射, 光滑且可表示的态射 $U\to\mathcal X\to X$, 由推论~\bookref{5.5.12} (或定理~\bookref{7.3.1}) 可知 $X$ 是代数空间. Gerbe 是光滑态射, 因而只需证明对每个来自概形的态射 $S\to X$, 纤维积 $\mathcal X\times_X S\to S$ 是 gerbe. 由构造, 存在 fppf 覆盖 $S'\to S$ 及提升复合 $S'\to S\to X$ 的态射 $a':S'\to\mathcal X$. 由于 gerbe 性质在目标上 fppf 局部, 将 $S$ 换为 $S'$ 后可假设 $S\to X$ 提升为映射 $a:S\to\mathcal X$. 此时有同构
\[
\Psi:\mathcal X\times_X S\longrightarrow B\underline{\operatorname{Aut}}_S(a).
\]
映射 $\Psi$ 将对象 $(b,f)\in(\mathcal X\times_X S)(T)$ 映到主 $\underline{\operatorname{Aut}}_S(a)$-丛 $\underline{\operatorname{Isom}}_T(f^*a,b)$. 请读者检查 $\Psi$ 是同构. 由于 $\underline{\operatorname{Aut}}_S(a)\to S$ 是 fppf 群概形, $B\underline{\operatorname{Aut}}_S(a)\to S$ 是 gerbe.

也可直接构造代数空间 $X$. 选取光滑呈示 $U\to\mathcal X$, 诱导光滑群胚 $R\rightrightarrows U$, 稳定子群胚概形 $S_U=R\times_{U\times U}U=I_{\mathcal X}\times_{\mathcal X}U$ 在 $U$ 上是 fppf. 存在 fppf 等价关系 $S_U\times_U R\rightrightarrows R$, 一个箭头由复合给出, 另一个为投影. 两次应用 fppf 等价关系的商的代数性 (推论~\bookref{7.3.6}), fppf 商 $R':=R/(S_U\times_U R)$ 是代数空间, 它诱导 fppf 等价关系 $R'\rightrightarrows U$, 且 $X$ 同构于 fppf 商 $U/R'$. 另见 \cite[推论 10.8]{LMB00} 和 \cite[Tags 06QD 与 06QJ]{SP}.
\end{proof}
\begin{exercise}\label{src:7.4.20}
证明每个 gerbe $\mathcal X\to X$, 其中 $\mathcal X$ 为代数叠, $X$ 为代数空间, 都是粗模空间 (定义~\bookref{5.2.1}).
\end{exercise}
\begin{exercise}\label{src:7.4.21}
若 $\mathcal X$ 是约化 Noether 代数叠, 证明存在稠密开子叠 $\mathcal U$, 它是 gerbe.

\textit{提示:} 使用一般平坦性 (\bookref{4.3.32}).
\end{exercise}
\begin{exercise}\label{src:7.4.22}
证明存在 $\mathbb A^1$ 上平凡带 $\mathbb Z/2$-带的 gerbe 的非平凡同构
\[
\alpha:B(\mathbb Z/2)\times(\mathbb A^1\setminus0)\longrightarrow B(\mathbb Z/2)\times(\mathbb A^1\setminus0),
\]
通过它可粘合得到 $\mathbb P^1$ 上非平凡的带 $\mathbb Z/2$-带的 gerbe.
\end{exercise}
\begin{exercise}\label{src:7.4.23}
设 $1\to K\to G\to Q\to1$ 是 $k$ 上仿射代数群的短正合序列, $K$ 交换. 证明 $BG\to BQ$ 是带 $K$-带的 gerbe, 它平凡当且仅当序列分裂.
\end{exercise}
\begin{exercise}\label{src:7.4.24}
假设 $\operatorname{char}(k)\neq2,3$. 回顾习题~\bookref{4.1.21}(b), 稳定椭圆曲线模叠有商描述 $\overline{\mathcal M}_{1,1}=\mathcal P(4,6):=[(\mathbb A^2\setminus0)/\mathbb G_m]$, 其中 $\mathbb G_m$ 以权 $4,6$ 作用.
\begin{enumerate}[(a)]
\item 证明 $j$-直线 $\pi:\mathcal M_{1,1}\to\mathbb A^1$ 在 $\mathbb A^1\setminus\{0,1728\}$ 上是平凡的带 $\mathbb Z/2$-带的 gerbe.

\textit{提示:} 通过 Weierstrass 方程
\[
y^2z+xyz=x^3-\frac{36}{t-1728}xz^2-\frac1{t-1728}z^3,
\]
构造 $\mathbb A_k^1\setminus\{0,1728\}$ 上的椭圆曲线族, 其中 $t$ 是 $\mathbb A^1$ 上的坐标, 判别式为 $\Delta=t^2/(t-1728)^3$. 见 \cite[命题 III.1.4(c)]{Sil09}.
\item 考虑由同态 $\mathbb G_m\to\mathbb G_m$, $t\mapsto t^2$, 诱导的加权射影叠之间的映射 $\overline{\mathcal M}_{1,1}=\mathcal P(4,6)\to\mathcal P(2,3)$ (稍后可将其视为刚性化的例子). 注意沿 $\mathcal P(2,3)\setminus\{0,1728,\infty\}$ 的限制是 (a) 中的 gerbe. 证明 $\overline{\mathcal M}_{1,1}\to\mathcal P(2,3)$ 非平凡.

\textit{提示:} 若它平凡, 证明存在挠线丛. 这与习题~\bookref{7.1.21}(a) 中 $\Pic(\overline{\mathcal M}_{1,1})=\mathbb Z$ 矛盾.
\item 证明 $\mathcal M_{1,1}\to\mathcal P(2,3)\setminus\infty$.

\textit{提示:} 若它平凡, 证明 $\mathcal M_{1,1}$ 有三个 $2$-挠线丛, 与习题~\bookref{7.1.21}(c) 中 $\Pic(\mathcal M_{1,1})=\mathbb Z/12$ 矛盾.
\end{enumerate}
\end{exercise}
\begin{exercise}\label{src:7.4.25}
\begin{enumerate}[(a)]
\item 证明代数空间上每个平展局部同构于 $B(\mathbb Z/2)$ 的 gerbe $\mathcal X$, 实际上带 $\mathbb Z/2$-带.
\item 给出代数空间上平展局部同构于 $B\mathbb G_m$, 但不带 $\mathbb G_m$-带的 gerbe 的例子.

\textit{提示:} 考虑 $\mathbb G_m$ 的一个形式的分类叠.
\end{enumerate}
\end{exercise}
\begin{exercise}[再论根 gerbe 与根叠]\label{src:7.4.26}
根 gerbe 与根叠已在例~\bookref{4.9.21} 和 \bookref{4.9.22} 中引入. 现在我们知道如何构造 $\mu_r$ 作用的商叠, 因而习题~\bookref{4.9.23} 可推广到 $r$ 在 $\Gamma(S,\mathcal O_S)$ 中未必可逆的情形. 给定概形 $X$ 及线丛 $L$, 证明 $X(\sqrt[r]{L})\to X$ 是带 $\mu_r$-带的 gerbe, 它平凡当且仅当 $L$ 有 $r$ 次根.
\end{exercise}
\subsection{附记: 非 Abel 带的一般知识}\label{src:7.4.4-subsection}
对景 $\mathcal S$ 上的群层 $G,H$, \newterm{外同构层}{sheaf of outer isomorphisms} $\underline{\operatorname{Out}}(G,H)$ 是预层
\[
T\longmapsto\operatorname{Out}(G|_T,H|_T):=\operatorname{Isom}(G|_T,H|_T)/H(T)
\]
的层化, 其中 $H(T)$ 通过共轭作用于 $\operatorname{Isom}(G|_T,H|_T)$. \emph{外同构} $G\to H$ 是 $\underline{\operatorname{Out}}(G,H)$ 在 $T$ 上的元素. 若 $H$ 为 Abel 群层, 则外同构与同构相同.
\begin{definition}[带的叠]\label{src:7.4.27}
设 $\mathcal S$ 是景. 定义 $\mathcal S$ 上的预叠 $\underline{\mathsf{Bands}}^{\mathrm{pre}}$, 其在 $T\in\mathcal S$ 上的对象是 $T$ 上的群层, 态射是外同构. 设 $\underline{\mathsf{Bands}}$ 是 $\underline{\mathsf{Bands}}^{\mathrm{pre}}$ 的叠化.
\end{definition}
\begin{remark}\label{src:7.4.28}
具体地, $\underline{\mathsf{Bands}}$ 在 $T\in\mathcal S$ 上的对象是如下数据: 覆盖 $\{T_i\to T\}$, 每个 $T_i$ 上的群层 $G_i$, 以及外同构 $\lambda_{ij}:G_i|_{T_{ij}}\to G_j|_{T_{ij}}$ (其中 $T_{ij}=T_i\times_T T_j$), 满足在 $T_{ijk}$ 上作为外同构的余圈条件 $\lambda_{jk}\circ\lambda_{ij}=\lambda_{ik}$. 使用相同覆盖的对象之间的态射
\[
\lambda:(\{T_i\to T\},G_i,\lambda_{ij})\longrightarrow(\{T_i\to T\},G_i',\lambda_{ij}')
\]
是外同构 $\phi_i:G_i\to G_i'$ 的数据, 满足 $\phi_i\circ\lambda_{ij}=\lambda_{ij}'\circ\phi_j$. 使用 $T$ 的不同覆盖的对象之间的态射, 由覆盖的共同加细及其限制之间的态射给出.
\end{remark}
\begin{definition}[带]\label{src:7.4.29}
设 $\mathcal S$ 是具有终对象 $X$ 的景. \newterm{带}{band} 是叠 $\underline{\mathsf{Bands}}$ 在 $X$ 上的对象 $(\{X_i\to X\},G_i,\lambda_{ij})$. 若每个 $G_i$ 是 Abel 群层, 则称该带为 \emph{Abel 带}.
\end{definition}
\begin{remark}\label{src:7.4.30}
$X$ 上的每个群层 $G$ 都可视为带. 但并非每个带都等价于层; 见习题~\bookref{7.4.38}. 事实上, 局部同构于 $G$ 的带的同构类集合可与 $\mathrm H^1(X,\underline{\operatorname{Out}}(G))$ 等同, 其中 $\underline{\operatorname{Out}}(G)$ 是 $G$ 的外自同构层; 见 \cite[推论 IV.1.1.7.3]{Gir71}.
\end{remark}
\begin{remark}\label{src:7.4.31}
$\underline{\mathsf{Bands}}$ 中 Abel 带构成的满子范畴, 等价于 Abel 群层的叠, 即它在 $T\in\mathcal S$ 上的对象是 $T$ 上的 Abel 群层.
\end{remark}
\begin{lemma}[Gerbe 的带]\label{src:7.4.32}
设 $\mathcal S$ 是具有终对象 $X$ 的景. $\mathcal S$ 上的每个 gerbe $\mathcal X$ 都诱导良定义的带 $B_{\mathcal X}$. Abel gerbe $\mathcal X$ 的带 $B_{\mathcal X}$ 总是群层.
\end{lemma}
\begin{proof}
可按习题~\bookref{7.4.15} 的方法证明.
\end{proof}
\begin{remark}\label{src:7.4.33}
若 $\mathcal S$ 没有终对象, 则 $\mathcal S$ 上的 gerbe $\mathcal X$ 诱导带的叠 $\underline{\mathsf{Bands}}\to\mathcal S$ 的截面 $\mathcal S\to\underline{\mathsf{Bands}}$.
\end{remark}
\begin{definition}[带 \texorpdfstring{$B$}{B}-带的 gerbe]\label{src:7.4.34}
设 $\mathcal S$ 是具有终对象 $X$ 的景, $B$ 是带. \emph{带 $B$-带的 gerbe} 是 gerbe $\mathcal X$ 配备带的同构 $B\xrightarrow{\sim}B_{\mathcal X}$. 带 $B$-带的 gerbe 的态射 $\mathcal X\to\mathcal Y$ 是叠态射, 使诱导的带映射 $B_{\mathcal X}\to B_{\mathcal Y}$ 与给定的从 $B$ 出发的同构交换.
\end{definition}
\begin{remark}\label{src:7.4.35}
若 $B$ 对应于 $\mathcal S$ 上的群层 $G$, 则带 $B$-带的 gerbe 的概念与带 $G$-带的 gerbe (定义~\bookref{7.4.10}) 一致.
\end{remark}
\begin{remark}[上同调刻画]\label{src:7.4.36}
若 $G$ 是景 $\mathcal S$ 上的 Abel 群层, 则上同调群 $\mathrm H^1(\mathcal S,G)$ 与 $\mathrm H^2(\mathcal S,G)$ 分别分类 $G$-挠子与带 $G$-带的 gerbe 的同构类 (见 \S\bookref{7.4.6-subsection}). 当 $G$ 非 Abel 时, 可将非 Abel 上同调带基点集合 $\mathrm H^1(\mathcal S,G)$ 与 $\mathrm H^2(\mathcal S,G)$ 定义为 $G$-挠子与带 $G$-带的 gerbe 的同构类, 这些集合也可用余圈等价地描述. 而且, 给定群层的短正合序列 $1\to G'\to G\to G''\to1$, 其中 $G'$ 为 Abel 群层且在 $G$ 中中心, 则诱导 $i\leq2$ 时第 $i$ 个上同调群的长正合序列. 见 \cite{Gir71}, \cite[144--145 页]{Mil80} 和 \cite[\S12]{Ols16}.
\end{remark}
\begin{exercise}[简单]\label{src:7.4.37}
证明存在不带 $\mathbb G_m$-带的 gerbe $\mathcal X$, 但对 $T$ 上每个对象 $x$, 仍有 (不具有函子性的) 同构 $\operatorname{Aut}_T(x)\cong\mathbb G_{m,T}$.
\end{exercise}
\begin{exercise}\label{src:7.4.38}
设 $G=\mathbb Z/5\rtimes\mathbb Z/4$, 其中 $\mathbb Z/4$ 作为 $\mathbb Z/5$ 的自同构群作用. 有群的正合序列
\[
1\longrightarrow D_5\longrightarrow G\longrightarrow\mathbb Z/2\longrightarrow1,
\]
其中 $D_5$ 是阶 $10$ 的二面体群; $D_5$ 的外自同构群为 $\mathbb Z/2$. 设 $K$ 是 $\mathbb Q$ 的虚二次扩张. 则 $P:=\Spec K\to\Spec\mathbb Q$ 是 $\mathbb Z/2$-挠子, 可考虑 $P$ 的提升 gerbe $\mathcal G_P$ (见习题~\bookref{7.4.45}). Gerbe $\mathcal X$ 平展局部同构于 $BD_5$. 证明 $\mathcal X$ 的带不同构于由 $\mathbb Q$ 上的群层诱导的带. (此例归功于 Kresch; 见 \cite[注记 3.2]{JL21}.)
\end{exercise}
\subsection{再论剩余 gerbe}\label{src:7.4.5-subsection}
给定代数叠 $\mathcal X$ 及 $x\in|\mathcal X|$, 回顾定义~\bookref{4.5.12}, $x$ 处的剩余 gerbe (若存在) 是约化且局部 Noether 的代数叠 $\mathcal G_x$, 配备单态射 $\mathcal G_x\hookrightarrow\mathcal X$, 使 $|\mathcal G_x|$ 是映到 $x$ 的点. 我们已经证明有限型点处的剩余 gerbe 存在且唯一 (命题~\bookref{4.5.16}).
\begin{proposition}[剩余 gerbe 的存在性 II]\label{src:7.4.39}
若 $\mathcal X$ 是 Noether 代数叠, $x\in|\mathcal X|$ 是点, 则剩余 gerbe $\mathcal G_x$ 存在, 且是某个域 $\kappa(x)$ 上的 gerbe, 称 $\kappa(x)$ 为 $x$ 的\emph{剩余域}.
\end{proposition}
\begin{proof}
将 $\mathcal X$ 换成闭包 $\overline{\{x\}}$, 可假设 $\mathcal X$ 约化, $x\in|\mathcal X|$ 稠密. 由习题~\bookref{7.4.21} 和定理~\bookref{5.5.1}, 将 $\mathcal X$ 换成开子叠后, 可假设 $\mathcal X$ 是概形 $X$ 上的 gerbe. 设 $y\in X$ 是 $x$ 的像, 则 $\mathcal G_x:=\mathcal X\times_X\kappa(y)$ 是 $\kappa(x):=\kappa(y)$ 上的 gerbe. 嵌入 $\mathcal G_x\to\mathcal X$ 是像为 $x$ 的单态射, $\mathcal G_x$ 是约化 Noether 子叠, 且 $|\mathcal G_x|$ 为单点集. 因而 $\mathcal G_x$ 是剩余 gerbe. 唯一性由命题~\bookref{4.5.16} 中相同的论证得到. 另见 \cite[定理 11.3]{LMB00}, \cite[定理 B.2]{Ryd11} 或 \cite[Tag 06UH]{SP}.
\end{proof}
若 $\mathcal X$ 是域 $\Bbbk$ 上拟分离有限型代数叠, $x\in\mathcal X(\Bbbk)$, 则 $\mathcal G_x=BG_x$ (命题~\bookref{4.5.16}). 更一般地, 有:
\begin{exercise}\label{src:7.4.40}
设 $\mathcal X$ 是 Noether 代数叠, $x\in|\mathcal X|$ 是有限型点.
\begin{enumerate}[(1)]
\item 对 $x$ 的任意代表 $\overline x:\Spec\Bbbk\to\mathcal X$, 有 Cartesian 图
\[
\begin{tikzcd}[ampersand replacement=\&]
BG_{\overline x}\arrow[r]\arrow[d]\&\mathcal G_x\arrow[r]\arrow[d]\&\mathcal X\\
\Spec\Bbbk\arrow[r]\arrow[ur,phantom,"{\square}"]\&\Spec\kappa(x).\&
\end{tikzcd}
\]
\item 若 $x$ 的稳定子光滑, 证明存在有限可分扩张 $\kappa(x)\to\Bbbk$ 及 $x$ 在 $\Bbbk$ 上的代表.
\end{enumerate}
\end{exercise}
\begin{exercise}\label{src:7.4.41}
设 $\mathcal X$ 是分离 Deligne--Mumford 叠, 粗模空间为 $\pi:\mathcal X\to X$. 证明 $x\in|\mathcal X|$ 的剩余 gerbe 是纤维 $\mathcal X_{\pi(x)}$ 的约化.
\end{exercise}
\begin{exercise}\label{src:7.4.42}
设 $C\subseteq\mathbb P_{\Bbbk}^2$ 是域 $\Bbbk$ 上的非分裂二次曲线, $\Bbbk\to\Bbbk'$ 是二次扩张, 使 $C\times_{\Bbbk}\Bbbk'\cong\mathbb P_{\Bbbk'}^1$. 设 $D\subseteq C$ 是次数 $6$ 的除子, $X\to\mathbb P_{\Bbbk'}^1$ 是在 $D\times_{\Bbbk}\Bbbk'$ 上分歧的二重覆盖. 证明 $[X]\in\mathcal M_2$ 的剩余 gerbe 非平凡, 剩余域为 $\Bbbk$. (事实上, 域上的每个 Deligne--Mumford gerbe 都可实现为某个亏格 $g$ 的曲线 $[C]\in\mathcal M_g$ 的剩余 gerbe; 见 \cite{BL24}.)
\end{exercise}
\subsection{上同调刻画}\label{src:7.4.6-subsection}
下列习题对概形 $S$ 的小 fppf 景 $S_{\mathrm{fppf}}$ 上的 Abel 层 $G$, 给出挠子与 gerbe 的上同调刻画. 若 $G\to S$ 是光滑, 交换且拟射影的群概形, 则 $\mathrm H^p(S_{\text{ét}},G)=\mathrm H^p(S_{\mathrm{fppf}},G)$ \cite[定理 3.9]{Mil80}. 对 $\mathbb G_m$, 有等同
\[
\Pic(S)=\mathrm H^1(S_{\mathrm{Zar}},\mathcal O_S^*)=\mathrm H^1(S_{\text{ét}},\mathbb G_m)=\mathrm H^1(S_{\mathrm{fppf}},\mathbb G_m)
\]
(Hilbert 第 $90$ 定理, \cite[命题 4.9]{Mil80}). 还应注意, 若 $S$ 是 $\mathbb C$ 上的光滑概形, $G$ 是有限 Abel 群, 则经典复上同调 $\mathrm H^p(S(\mathbb C),G)$ 与常值层的平展上同调 $\mathrm H^p(S_{\text{ét}},G)$ 一致 \cite[定理 3.12]{Mil80}.
作为额外挑战, 鼓励读者对任意景上的 Abel 层证明以下更一般的结论. 详细证明可查阅 \cite{Gir71} 和 \cite[\S12]{Ols16}.
\begin{exercise}[挠子]\label{src:7.4.43}
设 $S$ 是概形.
\begin{enumerate}[(a)]
\item 若 $G$ 是 $S_{\mathrm{fppf}}$ 上的 Abel 层, 证明 $\mathrm H^1(S_{\mathrm{fppf}},G)$ 与 $G$-挠子的同构类一一对应.

\textit{提示:} 模仿用 Čech 上同调证明概形 $X$ 满足 $\mathrm H^1(X,\mathcal O_X^*)=\Pic(X)$ 的论证.
\item 设 $0\to G'\to G\to G''\to0$ 是 $S_{\mathrm{fppf}}$ 上 Abel 层的正合序列, 对应的长正合序列为
\[
\begin{aligned}
0&\longrightarrow\mathrm H^0(S_{\mathrm{fppf}},G')\longrightarrow\mathrm H^0(S_{\mathrm{fppf}},G)\longrightarrow\mathrm H^0(S_{\mathrm{fppf}},G'')\\
&\xrightarrow{\delta}\mathrm H^1(S_{\mathrm{fppf}},G')\xrightarrow{\alpha}\mathrm H^1(S_{\mathrm{fppf}},G)\xrightarrow{\beta}\mathrm H^1(S_{\mathrm{fppf}},G'')\longrightarrow\cdots.
\end{aligned}
\]
证明在 (a) 的双射下, 边界映射 $\delta$ 将截面 $S\to G''$ 映到由纤维积 $G\times_{G''}S$ 定义的 $G'$-挠子. 还证明 $\alpha$ 将 $G'$-挠子 $P'$ 映到商 $P'\times_{G'}G:=(P'\times G)/G'$, 而 $\beta$ 将 $G$-挠子 $P$ 映到 $P\times_G G''$.
\end{enumerate}
\end{exercise}
\begin{exercise}[Gerbe]\label{src:7.4.44}
若 $G$ 是概形的 fppf 景 $S_{\mathrm{fppf}}$ 上的 Abel 层, 证明 $\mathrm H^2(S_{\mathrm{fppf}},G)$ 与带 $G$-带的 gerbe 的同构类一一对应.

\textit{提示:} 设 $0\to G\to I^0\xrightarrow{d^0}I^1\xrightarrow{d^1}I^2\xrightarrow{d^2}\cdots$ 是内射消解. 对上同调类 $\alpha\in\mathrm H^2(S_{\mathrm{fppf}},G)$, 如下定义 $S_{\mathrm{fppf}}$ 上的叠 $\mathcal G_\alpha$. 选取代表 $\alpha$ 的 $\tau\in\Gamma(S_{\mathrm{fppf}},I^2)$, 满足 $d^2(\tau)=0$. 定义 $\mathcal G_\alpha$ 为对 $(T,\sigma)$ 组成的范畴, 其中 $T\in S_{\mathrm{fppf}}$, $\sigma\in\Gamma(T,I^1)$ 满足 $d^1(\sigma)=\tau|_T$. 态射 $(T',\sigma')\to(T,\sigma)$ 是映射 $f:T'\to T$ 及元素 $\rho\in\Gamma(T',I^0)$ 的数据, 满足 $d^0(\rho)=\sigma'-f^*(\sigma)$. 证明 $\mathcal G_\alpha$ 是带 $G$-带的 gerbe, 且指派 $\alpha\mapsto\mathcal G_\alpha$ 给出所述双射. 另见 \cite[\S2.5]{dJ03}.
\end{exercise}
\begin{exercise}[续论 gerbe]\label{src:7.4.45}
设 $0\to G'\to G\to G''\to0$ 是概形的 fppf 景 $S_{\mathrm{fppf}}$ 上 Abel 层的正合序列, 对应的长正合序列为
\[
\begin{aligned}
\cdots&\longrightarrow\mathrm H^1(S_{\mathrm{fppf}},G'')\xrightarrow{\delta}\mathrm H^2(S_{\mathrm{fppf}},G')\\
&\xrightarrow{\alpha}\mathrm H^2(S_{\mathrm{fppf}},G)\xrightarrow{\beta}\mathrm H^2(S_{\mathrm{fppf}},G'')\longrightarrow\cdots.
\end{aligned}
\]
\begin{enumerate}[(a)]\setcounter{enumi}{1}
\item 证明在习题~\bookref{7.4.44} 的双射下, 边界映射 $\delta$ 将 $G''$-挠子 $P''\to S$ 映到\newterm{提升 gerbe}{gerbe of lifts} $\mathcal G_{P''}$, 其在 $S$-概形 $T$ 上的对象是对 $(P,\alpha)$, 由 $G$-挠子 $P\to T$ 及 $G''$-挠子的平凡化 $\alpha:P\times^G G''\cong P''\times_S T$ 组成.
\item 假设 $G'',G,G''$ 可由域 $\Bbbk$ 上的交换仿射代数群表示. 证明若 $P''\to S$ 是 $G''$-挠子, 则 $\mathcal G_{P''}$ 可同时与商叠 $[P''/G]$, 以及分类 $G''$-挠子 $P''$ 的映射 $S\to BG''$ 与 $BG\to BG''$ 的纤维积等同.
\item 在正合序列为 $1\to\boldsymbol\mu_r\to\mathbb G_m\to\mathbb G_m\to1$, 且 $P''$ 对应于线丛 $L''$ 的特殊情形, 证明 $\mathcal G_{P''}$ 同构于根叠 $X(\sqrt[r]{L})$.
\end{enumerate}
\end{exercise}
\begin{exercise}[群结构]\label{src:7.4.46}
证明 $\mathrm H^1(S_{\mathrm{fppf}},G)$ 与 $\mathrm H^2(S_{\mathrm{fppf}},G)$ 的群运算可如下作几何描述.
\begin{enumerate}[(a)]
\item 两个 $G$-挠子 $P_1,P_2$ 的乘积是\newterm{缩合积}{contracted product} $P_1\wedge^G P_2$, 定义为层商 $(P_1\times P_2)/G$, 其中 $h\cdot(p_1,p_2)=(h^{-1}p_1,hp_2)$, $G$-作用由 $g\cdot(p_1,p_2)=(gp_1,p_2)=(p_1,gp_2)$ 给出. $G$-挠子 $P$ 的逆是层 $P$, 配备取逆的 $G$-作用: $g\cdot p=g^{-1}p$.
\item 两个带 $G$-带的 gerbe $(\mathcal X_1,\psi_{1,x})$ 与 $(\mathcal X_2,\psi_{2,x})$ 的乘积是缩合积 $\mathcal X_1\wedge^G P_2$, 定义为乘积 $\mathcal X_1\times\mathcal X_2$ 沿带 $\psi_1,\psi_2$ 定义的子群 $(\psi_1,\psi_2):G|_{\mathcal X_1\times\mathcal X_2}\to I_{\mathcal X_1\times\mathcal X_2}$ 的刚性化 $(\mathcal X_1\times\mathcal X_2)\!\sslash G$ (见命题~\bookref{7.4.19}). 逆 $(\mathcal X,\psi_x)^{-1}=(\mathcal X,\psi_x^{-1})$ 将带取逆.
\end{enumerate}
\end{exercise}
\begin{remark}[\texorpdfstring{$\mathbb P^1$}{P1} 上带 \texorpdfstring{$\mu_n$}{mu n}-带的 gerbe]\label{src:7.4.47}
在代数闭域 $\Bbbk$ 上, $\mathbb P^1$ 上带 $\mu_n$-带的 gerbe 的同构类与 $\mathbb Z/n\mathbb Z$ 一一对应. 注意正合序列 $1\to\mu_n\to\mathbb G_m\xrightarrow{n}\mathbb G_m\to1$ 诱导上同调的正合序列
\[
\begin{aligned}
\mathrm H^1(\mathbb P^1_{\text{ét}},\mathbb G_m)&\xrightarrow{n}\mathrm H^1(\mathbb P^1_{\text{ét}},\mathbb G_m)\xrightarrow{\delta}\mathrm H^2(\mathbb P^1_{\text{ét}},\mu_n)\\
&\longrightarrow\mathrm H^2(\mathbb P^1_{\text{ét}},\mathbb G_m).
\end{aligned}
\]
由于 $\mathrm H^1(\mathbb P^1_{\text{ét}},\mathbb G_m)=\Pic(\mathbb P^1_{\text{ét}})=\mathbb Z$, 可用 $\mathrm H^2(\mathbb P^1_{\text{ét}},\mathbb G_m)=0$ 推出 $\mathrm H^2(\mathbb P^1_{\text{ét}},\mu_n)=\mathbb Z/n\mathbb Z$. 线丛 $\mathcal O(d)$ 在边界映射 $\delta$ 下的像等价于根叠 $\mathbb P^1(\sqrt[n]{\mathcal O(d)})$, 该 gerbe 平凡当且仅当 $n$ 整除 $d$.
\end{remark}
\begin{exercise}[群概形的形式]\label{src:7.4.48}
设 $G$ 是域 $\Bbbk$ 上的代数群. 若存在同构 $G_{\overline{\Bbbk}}\cong H_{\overline{\Bbbk}}$, 则称群概形 $H\to\Spec\Bbbk$ 为 $G$ 的\newterm{形式}{form}. 称 $G$ 为 $G$ 的\emph{平凡形式}.
\begin{enumerate}[(a)]
\item 证明 $\mathbb R$ 上的代数群 $H=\Spec\mathbb R[x,y]/(x^2+y^2-1)$, 群结构由嵌入 $H\subseteq\mathrm{SL}_2$,
\[
(x,y)\longmapsto\begin{pmatrix}x&y\\-y&x\end{pmatrix},
\]
诱导, 是 $\mathbb G_{m,\mathbb R}$ 的非平凡形式.
\item 假设 $\operatorname{char}(\Bbbk)\neq2$. 回顾习题~\bookref{B.1.60} 中关于 $n$ 维向量空间 $V$ 上非退化二次型 $q$ 定义的正交群 $O(q)$. 证明每个 $O(q)$ 都是正交矩阵子群 $O_n\subseteq\mathrm{GL}_n$ 的形式.
\item 若 $G$ 光滑且交换, 证明 $G$ 的形式由 $\mathrm H^1((\Sch/\Bbbk)_{\text{ét}},\operatorname{Aut}(G))$ 分类.
\end{enumerate}
\end{exercise}
\begin{exercise}[\texorpdfstring{$\mathbb G_m$}{Gm}-gerbe 与扭曲层]\label{src:7.4.49}
设 $\mathcal X\to X$ 是代数空间 $X$ 上的 $\mathbb G_m$-gerbe. 若对每个域值点 $x:\Spec\Bbbk\to\mathcal X$, $E$ 沿 $B\mathbb G_m=BG_x\to\mathcal X$ 的拉回所对应的 $\mathbb G_m$-表示分解为权 $1$ 的一维表示的直和, 则称 $\mathcal X$ 上的凝聚层 $F$ 为 \newterm{$1$-扭曲的}{$1$-twisted}. 证明 $\mathbb G_m$-gerbe $\mathcal X\to X$ 平凡, 当且仅当 $\mathcal X$ 上存在 $1$-扭曲线丛.
\end{exercise}
\begin{exercise}[Azumaya 代数]\label{src:7.4.50}
Noether 概形 $X$ 上\newterm{秩 $r^2$ 的 Azumaya 代数}{Azumaya algebra of rank $r^2$} 是 (可能非交换的) 结合 $\mathcal O_X$-代数 $A$, 作为 $\mathcal O_X$-模凝聚, 且存在平展覆盖 $X'\to X$, 使 $A\otimes_{\mathcal O_X}\mathcal O_{X'}$ 同构于矩阵代数 $\mathrm{Mat}_{r\times r}(\mathcal O_X)$. 若 $A$ 同构于 $\mathrm{Mat}_{r\times r}(\mathcal O_X)$, 则称它\emph{平凡}. 由习题~\bookref{B.1.70}, Azumaya 代数与主 $\mathrm{PGL}_r$-丛, 以及相对维数 $r$ 的 Brauer--Severi 概形一一对应. 设 $A$ 是 Noether 概形 $X$ 上秩 $r^2$ 的 Azumaya 代数.
\begin{enumerate}[(a)]
\item 定义 $A$ 的\emph{平凡化 gerbe} 为 $(\Sch/X)_{\text{ét}}$ 上的叠 $\mathcal G_A$, 其在 $X$-概形 $T$ 上的对象是对 $(E,\alpha)$, 由 $T$ 上秩 $r$ 的向量丛 $E$ 及平凡化 $\alpha:\operatorname{End}_{\mathcal O_X}(E)\xrightarrow{\sim}A\otimes_{\mathcal O_X}\mathcal O_T$ 组成. $\mathcal G_A(T)$ 中的态射是与平凡化相容的向量丛同构. 证明 $\mathcal G_A\to X$ 是带 $\mathbb G_m$-带的 gerbe.
\item 设 $P_A$ 是对应于 $A$ 的主 $\mathrm{PGL}_r$-丛. 将 $\mathcal G_A$ 与习题~\bookref{7.4.45} 中相对于 $\mathrm{PGL}_r$-挠子 $P_A$ 及满射 $\mathrm{GL}_r\to\mathrm{PGL}_r$ 定义的提升 gerbe $\mathcal G_{P_A}$ 等同.
\item 用四元数构造 $\Spec\mathbb R$ 上非平凡的 $\mathbb G_m$-gerbe.
\end{enumerate}
\end{exercise}
\begin{remark}[Brauer 群]\label{src:7.4.51}
若 Noether 概形 $X$ 上两个 Azumaya 代数 $A,A'$ 满足: 存在 $X$ 上的向量丛 $E,E'$, 使 $A\otimes_{\mathcal O_X}\operatorname{End}_{\mathcal O_X}(E)\cong A\otimes_{\mathcal O_X}\operatorname{End}_{\mathcal O_X}(A')$, 则称它们\emph{相似}. 这定义一个等价关系, $X$ 的 \newterm{Brauer 群}{Brauer group} 是 Azumaya 代数的相似类集合 $\operatorname{Br}(X)$. 集合 $\operatorname{Br}(X)$ 在运算 $[A]\cdot[A']=[A\otimes A']$ 和 $[A]^{-1}=[A^{\mathrm{op}}]$ 下成为群 (其中 $A^{\mathrm{op}}$ 是反代数, 元素及加法与 $A$ 相同, 乘法反向: $a\cdot_{A^{\mathrm{op}}}b=b\cdot_A a$).

$X_{\text{ét}}$ 上层的正合序列 $1\to\mathbb G_{m,X}\to\mathrm{GL}_{r,X}\to\mathrm{PGL}_{r,X}\to1$ 诱导边界映射
\[
\delta:\mathrm H^1(X_{\text{ét}},\mathrm{PGL}_r)\longrightarrow\mathrm H^2(X_{\text{ét}},\mathbb G_m),\qquad P_A\longmapsto\mathcal G_A,
\]
其中 $P_A$ 是对应于 $A$ 的主 $\mathrm{PGL}_r$-丛, $\mathcal G_A$ 是 $A$ 的平凡化 gerbe. 不难检查 $[\mathcal G_A]$ 是 $r$-挠元素, 且 $A$ 平凡当且仅当 $\mathcal G_A$ 平凡. 因而存在到\newterm{上同调 Brauer 群}{cohomological Brauer group} $\operatorname{Br}'(X)$ 的单射
\[
\operatorname{Br}(X)\hookrightarrow\operatorname{Br}'(X):=\mathrm H^2(X_{\text{ét}},\mathbb G_m)_{\mathrm{tors}},\qquad A\longmapsto\mathcal G_A.
\]
进一步背景见 \cite{Gro68} 和 \cite[\S IV.2]{Mil80}. Grothendieck 问过 $\operatorname{Br}(X)\hookrightarrow\operatorname{Br}'(X)$ 是否满射. 某些情形已知. 最强的结果归功于 Gabber: 若 $X$ 有丰沛线丛, 则 $\operatorname{Br}(X)=\operatorname{Br}'(X)$ (见 \cite{dJ03}). 但一般情形仍未解决, 即便对域上的光滑分离概形也是如此.
\end{remark}
\begin{exercise}\label{src:7.4.52}
设 $X$ 是 Noether 概形, $\mathcal X\to X$ 是对应于上同调类 $[\mathcal X]\in\mathrm H^2(X_{\text{ét}},\mathbb G_m)$ 的带 $\mathbb G_m$-带的 gerbe.
\begin{enumerate}[(a)]
\item 证明下列条件等价:
\begin{enumerate}[(i)]
\item 存在 $X$ 上的 Azumaya 代数 $A$, 使 $\mathcal X\cong\mathcal G_A$, 即 $[\mathcal X]$ 属于 $\operatorname{Br}'(X)\to\operatorname{Br}(X)$ 的像.
\item $\mathcal X$ 是全局商叠.
\item $\mathcal X$ 上存在 $1$-扭曲向量丛 $E$ (见习题~\bookref{7.4.49}).
\end{enumerate}
\item 设 $X$ 是 $\mathbb C$ 上的正规分离曲面, $\mathrm H^2(X,\mathbb G_m)$ 含有非挠元素 $\alpha$; 一个例子见 \cite[II.1.11.b]{Gro68}. 推出对应于 $\alpha$ 的带 $\mathbb G_m$-带的 gerbe 不是全局商叠.
\item 设 $Y=\Spec\mathbb C[x,y,z]/(xy-z^2)$. 证明 $(Y\setminus0)\times B(\mathbb Z/2)$ 存在非平凡对合 $\alpha$, 使沿 $\alpha$ 粘合 $Y$ 上平凡的带 $\mathbb Z/2$-带的 gerbe 所得的叠 $\mathcal X$, 是非分离并 $Y\cup_{Y\setminus0}Y$ 上的带 $\mathbb Z/2$-带的 gerbe, 且 $\mathcal X$ 不是全局商叠.
\end{enumerate}
另见 \cite{EHKV01}.
\end{exercise}
\subsection{刚性化}\label{src:7.4.7-subsection}
给定代数叠 $\mathcal X$, 惯性叠 $I_{\mathcal X}$ 可视为 $\mathcal X$ 的大平展景 $(\Sch/\mathcal X)_{\text{ét}}$ 上的群代数空间. 作为群函子, $I_{\mathcal X}$ 将对象 $a\in\mathcal X(S)$ 映到群 $\operatorname{Aut}_S(a)$, 将 $S'\to S$ 上的态射 $\alpha:a'\to a$ 映到自然拉回映射 $\alpha^*:\operatorname{Aut}_S(a)\to\operatorname{Aut}_{S'}(a')$ (见 \bookeqref{4.2.11}). 给定 $a:S\to\mathcal X$, 有 $S$ 上群代数空间的典范同构 $I_{\mathcal X}\times_{\mathcal X}S\cong\underline{\operatorname{Aut}}_S(a)$.
假设 $\mathcal H\subseteq I_{\mathcal X}$ 是 $\mathcal X$ 上的闭子群空间, 且 $\mathcal H\to\mathcal X$ 为 fppf. 等价地, 要求对每个 $a\in\mathcal X(S)$, 存在 $S$ 上闭 fppf 子群概形 $\mathcal H_a\subseteq\underline{\operatorname{Aut}}_S(a)$, 使若 $a'\to a$ 是 $S'\to S$ 上的态射, 则典范同构 $\underline{\operatorname{Aut}}_{S'}(a')\cong\underline{\operatorname{Aut}}_S(a)\times_S S'$ 限制为同构 $\mathcal H_{a'}\cong\mathcal H_a\times_S S'$. 若 $\alpha:a\xrightarrow{\sim}a$ 是恒等态射上的自同构, 则典范同构 $\alpha^*:\underline{\operatorname{Aut}}_S(a)\to\underline{\operatorname{Aut}}_S(a)$ 是由 $\alpha$ 给出的共轭. 特别地, $\mathcal H_a\subseteq\underline{\operatorname{Aut}}_S(a)$ 是正规群概形.
\begin{definition}[刚性化]\label{src:7.4.53}
设 $\mathcal X$ 是代数叠, $\mathcal H\subseteq I_{\mathcal X}$ 是 $\mathcal X$ 上的 fppf 闭子群概形. \newterm{刚性化}{rigidification} $\mathcal X\!\sslash\mathcal H$ 定义为如下预叠在 $\Sch_{\mathrm{fppf}}$ 中的叠化: 对象与 $\mathcal X$ 相同, $f:S\to T$ 上从 $a\in\mathcal X(S)$ 到 $b\in\mathcal X(T)$ 的态射集合定义为
\[
\operatorname{Mor}(a,b)=\operatorname{Mor}_{\mathcal X(S)}(a,f^*b)/\mathcal H(T).
\]
\end{definition}
可将子群 $\mathcal H$ 视为给出 $B\mathcal H$ 在 $\mathcal X$ 上的作用, 将刚性化 $\mathcal X\!\sslash\mathcal H$ 视为商 $\mathcal X/B\mathcal H$. 应用中, 闭子群 $\mathcal H\subseteq I_{\mathcal X}$ 往往来自 fppf 群概形 $H\to S$ 的拉回, 即 $\mathcal H=H\times_S\mathcal X$. 此时通常写成 $\mathcal X\!\sslash H:=\mathcal X\!\sslash\mathcal H$.
\begin{example}\label{src:7.4.54}
若 $I_{\mathcal X}\to\mathcal X$ 是 fppf, 可取 $\mathcal H=I_{\mathcal X}$, 刚性化 $\mathcal X\!\sslash I_{\mathcal X}$ 是命题~\bookref{7.4.19} 中构造的代数空间 $X$.
\end{example}
\begin{proposition}\label{src:7.4.55}
设 $\mathcal X$ 是代数叠, $\mathcal H\subseteq I_{\mathcal X}$ 是 $\mathcal X$ 上的 fppf 闭子群概形. 刚性化 $\mathcal X\!\sslash\mathcal H$ 是代数叠, 且满足
\begin{enumerate}[(1)]
\item 自然态射 $\pi:\mathcal X\to\mathcal X\!\sslash\mathcal H$ 是 gerbe.
\item 对每个对象 $a\in\mathcal X(S)$, 自然映射 $\operatorname{Aut}_S(a)\to\operatorname{Aut}_S(\pi(a))$ 满射, 核为 $\mathcal H(S)$.
\item 态射 $f:\mathcal X\to\mathcal Y$ 唯一地经 $\mathcal X\!\sslash\mathcal H$ 分解, 当且仅当对每个对象 $a\in\mathcal X(S)$,
\[
\mathcal H(S)\subseteq\ker\bigl(\operatorname{Aut}_{\mathcal X(S)}(a)\to\operatorname{Aut}_{\mathcal Y(S)}(f(a))\bigr).
\]
\item 若 $\mathcal H$ 是交换群概形, 则 $\mathcal H$ 下降为 fppf 群概形 $H\to X$, 使 $\mathcal X\to X$ 是带 $H$-带的 gerbe. 若 $\mathcal H=H\times_S\mathcal X$ 还是交换 fppf 群概形 $H\to S$ 的拉回, 则 $\mathcal X\to X$ 是带 $H$-带的 gerbe.
\end{enumerate}
\end{proposition}
\begin{proof}
为证明 $\mathcal X$ 是代数的, 只需证明 $\pi:\mathcal X\to\mathcal X\!\sslash\mathcal H$ 是光滑可表示态射. 给定 $g:S\to\mathcal X\!\sslash\mathcal H$, 由 $\mathcal X\!\sslash\mathcal H$ 的叠化定义, 存在 fppf 覆盖 $S'\to S$, 使 $S'\to S\to\mathcal X\!\sslash\mathcal H$ 提升为映射 $a':S'\to\mathcal X$. 将 $S$ 换成 $S'$ 后, 可假设 $g:S\to\mathcal X\!\sslash\mathcal H$ 提升为态射 $a:S\to\mathcal X$. 如命题~\bookref{7.4.19} 的证明, 可证明存在同构
\[
\Psi:\mathcal X\times_{\mathcal X\sslash\mathcal H}S\longrightarrow B\mathcal H_a.
\]
细节留给读者. 另见 \cite[定理 5.1.5]{ACV03}, \cite[\S C]{AGV08}, \cite[\S5]{Rom05} 和 \cite[\S A]{AOV08}.
\end{proof}
\begin{example}\label{src:7.4.56}
每个稳定椭圆曲线族都有典范的超椭圆对合. 这给出域上稳定椭圆曲线叠 $\overline{\mathcal M}_{1,1}\cong\mathcal P(4,6)$ 的惯性叠 $I_{\overline{\mathcal M}_{1,1}}$ 的子群 $\mu_2:=\mu_{2,\overline{\mathcal M}_{1,1}}$ (照例假设 $\operatorname{char}(\Bbbk)\neq2,3$). 刚性化 $\overline{\mathcal M}_{1,1}\to\overline{\mathcal M}_{1,1}\!\sslash\mu_2$ 可与习题~\bookref{7.4.24} 中讨论的非平凡 $\mu_2$-gerbe $\mathcal P(4,6)\to\mathcal P(2,3)$ 等同.
\end{example}
\begin{example}[\texorpdfstring{$\mathcal B\!un(C)$}{Bun(C)} 的刚性化]\label{src:7.4.57}
层的模叠给出有趣的刚性化例子, 因为层上有典范的伸缩 $\mathbb G_m$-作用. 例如, 考虑代数闭域 $\Bbbk$ 上固定的光滑, 连通射影曲线 $C$ 上的向量丛模叠 $\mathcal B\!un_{r,d}(C)$. 对 $\Bbbk$-概形 $S$ 及 $C\times S$ 上任意向量丛 $\mathcal E$, 有 $S$ 上群概形的典范闭浸入 $i_{\mathcal E}:\mathbb G_{m,S}\to\underline{\operatorname{Aut}}(\mathcal E)$. 因而 $\mathbb G_m:=\mathbb G_{m,\mathcal B\!un_{r,d}}\subseteq I_{\mathcal B\!un_{r,d}}$ 是惯性叠的闭 fppf 群概形, 可以构造刚性化
\[
\mathcal B\!un(C)\!\sslash\mathbb G_m.
\]
在单向量丛 (即满足 $\operatorname{Aut}(E)=\Bbbk^*$ 的向量丛 $E$) 的开子叠 $\mathcal B\!un^{\mathrm{simple}}(C)$ 上, 刚性化 $\mathcal B\!un^{\mathrm{simple}}(C)\!\sslash\mathbb G_m$ 是代数空间, 因而 $\mathcal B\!un^{\mathrm{simple}}(C)$ 是其粗模空间上的 $\mathbb G_m$-gerbe.
\end{example}
\begin{exercise}\label{src:7.4.58}
\begin{enumerate}[(a)]
\item 若 $H$ 是 $S$ 上的交换 fppf 群概形, 证明 $BH\!\sslash H\cong S$. 更一般地, 若 $\mathcal X\to X$ 是带 $H$-带的 gerbe, 证明 $X\cong\mathcal X\!\sslash H$.
\item 设 $G\to S$ 是 fppf 群概形, 作用于 $S$-概形 $U$. 假设 $H\subseteq G$ 是在 $U$ 上平凡作用的中心交换 fppf 子群概形. 证明 $[U/G]\!\sslash H\cong[U/(G/H)]$.
\end{enumerate}
\end{exercise}
\begin{exercise}\label{src:7.4.59}
设 $\mathcal X$ 是概形 $S$ 上光滑, 整且分离的 Deligne--Mumford 叠. 设 $\Spec K\to\mathcal X$ 是一般点的代表. 证明惯性的一般纤维 $I_{\mathcal X}\times_{\mathcal X}K$ 的闭包 $\mathcal H\subseteq I_{\mathcal X}$ 是闭平展子群概形, 且刚性化 $\mathcal X\!\sslash\mathcal H$ 是 $S$ 上光滑, 整且分离的 Deligne--Mumford 叠, 一般惯性平凡.
\end{exercise}
\begin{exercise}\label{src:7.4.60}
设 $\mathcal X$ 是概形 $S$ 上的代数叠, $H\to S$ 是 fppf 群概形, $H\times_S\mathcal X\subseteq I_{\mathcal X}$ 是闭子群概形. 证明刚性化 $\mathcal X\!\sslash H$ 有如下模解释: 在概形 $S$ 上的对象是对 $(\mathcal G,f)$, 其中 $\mathcal G\to S$ 是带 $H$-带的 gerbe, $f:\mathcal G\to\mathcal X$ 是 $H$-等变态射.
\end{exercise}
\begin{exercise}\label{src:7.4.61}
对 $g\geq2$, 设 $\mathcal H_g\subseteq\mathcal M_g$ 是分类超椭圆曲线的闭子叠. 证明沿超椭圆对合的刚性化 $\mathcal H_g\to\mathcal Y$ 是非平凡的带 $\mathbb Z/2$-带的 gerbe.
\end{exercise}
\subsection{Picard 叠与空间}\label{src:7.4.8-subsection}
若 $X\to S$ 是 Noether 概形的固有平坦态射, 定义 \newterm{Picard 叠}{Picard stack}\footnote{文献中, ``Picard 叠''也指叠的 $2$-范畴中的群对象. 这是 Deligne 首先在 \cite[XVIII.1.4]{SGA4} 中引入的有用概念.}
\[
\underline{\mathcal P\!ic}_{X/S}
\]
为 $(\Sch/S)_{\text{ét}}$ 上的叠, 其在 $S$-概形 $T$ 上的对象是 $X_T=X\times_S T$ 上的线丛, 态射是线丛的同构. 它是 $\underline{\mathrm{Coh}}(X/S)$ 的开子叠, 后者是 $S$ 上局部有限型的代数叠; 当 $X\to S$ 强射影时, 这由使用 Quot 概形的显式呈示得到 (习题~\bookref{4.1.25}), 一般情形则是 Artin 公理 (定理~\bookref{C.7.7}) 的结果. 因而 $\underline{\mathcal P\!ic}_{X/S}$ 也是 $S$ 上局部有限型的代数叠.
另一方面, Picard 函子有几个候选:
\begin{enumerate}[(1)]
\item \newterm{朴素 Picard 函子}{naive Picard functor} (或\emph{绝对 Picard 函子}) 是
\[
\underline{\Pic}_{X/S}^{\mathrm{naive}}:\Sch/S\longrightarrow\mathsf{Gps},\qquad(T\to S)\longmapsto\Pic(X_T).
\]
\item \newterm{Picard 函子}{Picard functor}
\[
\underline{\Pic}_{X/S}:\Sch/S\longrightarrow\mathsf{Gps}
\]
是 $\underline{\Pic}_{X/S}^{\mathrm{naive}}$ 的 fppf 层化.
\item \newterm{相对 Picard 函子}{relative Picard functor} 是
\[
\underline{\Pic}_{X/S}^{\mathrm{rel}}:\Sch/S\longrightarrow\mathsf{Gps},\qquad(T\to S)\longmapsto\Pic(X_T)/\Pic(T),
\]
其在 $S$-概形 $T$ 上的对象是 $X_T$ 上的线丛 $L$; 若存在 $M\in\Pic(T)$ 使 $L\cong L'\otimes f_T^*M$, 则将 $X_T$ 上的两个线丛 $L,L'$ 等同.
\item 在对任意映射 $T\to S$, $\mathcal O_T\xrightarrow{\sim}f_{T,*}\mathcal O_{X_T}$ 都是同构的假设下, 可以定义 Picard 叠的刚性化
\[
\underline{\mathcal P\!ic}_{X/S}\!\sslash\mathbb G_m.
\]
这一假设表明, 对 $X_T$ 上线丛 $L$, 有典范同构
\[
\mathbb G_m(T)=\Gamma(T,\mathcal O_T)^*\xrightarrow{\sim}\Gamma(X_T,\mathcal O_{X_T})^*=\operatorname{Aut}(L),
\]
从而惯性叠 $I_{\underline{\mathcal P\!ic}_{X/S}}$ 同构于 $\underline{\mathcal P\!ic}_{X/S}$ 上的 fppf 群概形 $\mathbb G_m:=\mathbb G_{m,\underline{\mathcal P\!ic}_{X/S}}$.
\end{enumerate}
朴素 Picard 函子 $\underline{\Pic}_{X/S}^{\mathrm{naive}}$ 不是大 Zariski 拓扑中的层, 因为余圈条件未必满足, 因而它不可表示. 另一方面, Picard 函子 $\underline{\Pic}_{X/S}$ 与刚性化 $\underline{\mathcal P\!ic}_{X/S}\!\sslash\mathbb G_m$ 按定义是大 fppf 拓扑中的层. 在相当温和的假设下, $\underline{\Pic}_{X/S}^{\mathrm{rel}}$ 也是层, 这可能令人惊讶.
\begin{proposition}\label{src:7.4.62}
设 $f:X\to S$ 是 Noether 概形的固有平坦态射, 使 $\mathcal O_S\xrightarrow{\sim}f_*\mathcal O_X$ 是同构, 且基变换后仍成立, 即对每个映射 $g:T\to S$, 映射 $\mathcal O_T\xrightarrow{\sim}f_{T,*}\mathcal O_{X_T}$ 是同构. 则
\begin{enumerate}[(1)]
\item $\underline{\Pic}_{X/S}$ 可由 $S$ 上局部有限型的代数空间表示.
\item Picard 叠的映射 $\underline{\mathcal P\!ic}_{X/S}\to\underline{\Pic}_{X/S}$ 是带 $\mathbb G_m$-带的 gerbe, 且 $\underline{\Pic}_{X/S}\cong\underline{\mathcal P\!ic}_{X/S}\!\sslash\mathbb G_m$.
\item 若还存在截面 $s:S\to X$, 则 $\underline{\Pic}_{X/S}\cong\underline{\Pic}_{X/S}^{\mathrm{rel}}$.
\end{enumerate}
若 $f$ 的几何纤维还整, 则 $\underline{\Pic}_{X/S}$ 在 $S$ 上分离.
\end{proposition}
\begin{remark}\label{src:7.4.63}
若 $f:X\to S$ 是具有几何连通且约化纤维的固有平坦态射, 则由引理~\bookref{A.6.13}, $\mathcal O_S\xrightarrow{\sim}f_*\mathcal O_X$ 是同构 (且基变换后仍如此). Grothendieck 在 \cite[定理 3.1]{FGAV} 中证明, 当 $X\to S$ 射影时 $\underline{\Pic}_{X/S}$ 是概形. 其他阐述及各种推广见 \cite[\S20--21]{Mum66a}, \cite{AK80} 和 \cite[\S9.4]{Kle05}. 上述代数空间可表示性首先由 Artin 建立 \cite[定理 7.3]{Art69b}, 在稍弱的假设下也成立, 即 $f$ 在维数 $0$ 上上同调平坦: $f_*\mathcal O_X$ 的形成与基变换交换.
\end{remark}
\begin{proof}
如上所述, $\mathcal O_S\xrightarrow{\sim}f_*\mathcal O_X$ 在基变换后仍成立这一条件, 表明 $\underline{\mathcal P\!ic}_{X/S}$ 的惯性叠同构于 fppf 群概形 $\mathbb G_m:=\mathbb G_{m,\underline{\mathcal P\!ic}_{X/S}}$. 因而命题~\bookref{7.4.19} 表明刚性化 $\underline{\mathcal P\!ic}_{X/S}\!\sslash\mathbb G_m$ 是 $S$ 上局部有限型的代数空间. 而且, 由刚性化的定义, 它可与 Picard 函子 $\underline{\Pic}_{X/S}$ 等同. 这同时给出 (1) 与 (2). 当纤维几何整时, $\underline{\Pic}_{X/S}$ 在 $S$ 上的分离性由命题~\bookref{A.6.18} 和注记~\bookref{A.6.20} 得到.

为将 $\underline{\Pic}_{X/S}^{\mathrm{rel}}$ 与 $\underline{\Pic}_{X/S}$ 等同, 只需证明前者是大 fppf 拓扑中的层. 为此, 将 $\underline{\Pic}_{X/S}^{\mathrm{rel}}$ 与预叠 $\mathcal P^{s\text{-rig}}$ 等同更为方便, 称它为\emph{由截面 $s$ 给出的刚性化}, 其在 $S$-概形 $T$ 上的纤维范畴是
\[
\mathcal P^{s\text{-rig}}(T)=\{(L,\alpha)\mid L\in\Pic(X_T),\ \alpha:\mathcal O_T\xrightarrow{\sim}s_T^*L\},
\]
其中态射 $(L,\alpha)\to(L',\alpha')$ 是线丛同构 $\beta:L\to L'$, 满足 $\alpha'=s_T^*\beta\circ\alpha$. 考虑 $\mathcal P^{s\text{-rig}}$ 的优点是, 直接应用拟凝聚层的 fppf 下降即可检查它是大 fppf 拓扑 $(\Sch/S)_{\mathrm{fppf}}$ 上的叠. 因而只需验证自然映射
\begin{equation}\label{src:7.4.64}
\mathcal P^{s\text{-rig}}\longrightarrow\underline{\Pic}_{X/S}^{\mathrm{rel}},\qquad(L,\alpha)\longmapsto L
\end{equation}
是等价.

先检查 \bookeqref{7.4.64} 忠实, 特别地 $\mathcal P^{s\text{-rig}}$ 等价于函子. 必须证明若 $\beta:L\xrightarrow{\sim}L$ 是自同构且 $s_T^*\beta=\mathrm{id}_{s_T^*L}$, 则 $\beta=\mathrm{id}_L$. 由于 $\mathcal O_T\to f_{T,*}\mathcal O_{X_T}$ 是同构, 拉回映射 $f_T^*:\mathrm H^0(T,\mathcal O_T)\to\mathrm H^0(X_T,\mathcal O_{X_T})$ 是同构; $s_T$ 为截面, 故 $s_T^*$ 是 $f_T^*$ 的逆. 复合
\[
\begin{aligned}
&s_T^*:\operatorname{Hom}_{\mathcal O_{X_T}}(L,L)\cong\mathrm H^0(X_T,\mathcal O_{X_T})\\
&\hspace{4em}\xrightarrow{s_T^*}\mathrm H^0(T,\mathcal O_T)\cong\operatorname{Hom}_{\mathcal O_T}(s_T^*L,s_T^*L)
\end{aligned}
\]
是群同构, 因而 $\beta=\mathrm{id}_L$. 为证明 \bookeqref{7.4.64} 是满的 (即同构类函子上的诱导映射单射), 设 $(L,\alpha)\in\mathcal P^{s\text{-rig}}(T)$, 存在 $T$ 上的线丛 $M$ 及同构 $\beta:L\xrightarrow{\sim}f_T^*M$. 同构
\[
\begin{aligned}
\gamma:\mathcal O_{X_T}=f_T^*\mathcal O_T&\xrightarrow{f_T^*\alpha}f_T^*s_T^*L\\
&\xrightarrow{f_T^*s_T^*\beta}f_T^*s_T^*f_T^*M=f_T^*M
\end{aligned}
\]
满足 $s_T^*\gamma=s_T^*\beta\circ\alpha$, 表明 $(L,\alpha)$ 同构于 $(\mathcal O_{X_T},\mathrm{id})\in\mathcal P^{s\text{-rig}}(T)$. 最后, 为证明函子 \bookeqref{7.4.64} 本质满, 设 $L\in\Pic(X_T)$, 定义 $L'=L\otimes(f_T^*s_T^*L)^\vee$. $L,L'$ 在 $\underline{\Pic}_{X/S}^{\mathrm{rel}}(T)$ 中的像相同, 且 $s_T^*L'\cong s_T^*L\otimes(s_T^*f_T^*s_T^*L)^\vee\cong\mathcal O_T$ 定义同构 $\alpha':\mathcal O_T\xrightarrow{\sim}s_T^*L'$, 使 $L\in\underline{\Pic}_{X/S}^{\mathrm{rel}}(T)$ 是 $(L',\alpha')\in\mathcal P^{s\text{-rig}}(T)$ 的像.
\end{proof}
在代数闭域 $\Bbbk$ 上, 验证 Picard 函子 $\underline{\Pic}_X:=\underline{\Pic}_{X/\Bbbk}$ 是概形非常容易.
\begin{theorem}\label{src:7.4.65}
设 $X$ 是代数闭域 $\Bbbk$ 上的固有整概形.
\begin{enumerate}[(1)]
\item $\underline{\Pic}_X$ 是 $\Bbbk$ 上局部有限型的群概形, 特别地, 单位元的连通分支 $\underline{\Pic}_X^0$ 拟射影.
\item $\underline{\Pic}_X\cong\underline{\Pic}_X^{\mathrm{rel}}$, 且 $\underline{\mathcal P\!ic}_X\to\underline{\Pic}_X$ 是带 $\mathbb G_m$-带的 gerbe.
\item 若 $X$ 光滑, 则 $\underline{\Pic}_X^0$ 射影.
\item 若 $\operatorname{char}(\Bbbk)=0$, 则 $\underline{\Pic}_X$ 光滑, 维数为 $h^0(X,\mathcal O_X)$. 特别地, $\underline{\Pic}_X^0$ 是 Abel 簇.
\end{enumerate}
\end{theorem}
\begin{proof}
$\Bbbk$ 代数闭且 $X$ 整, 因而结构映射 $f:X\to\Spec\Bbbk$ 有截面, 对任意 $\Bbbk$-概形 $T$, $\mathcal O_T\xrightarrow{\sim}f_{T,*}\mathcal O_{X_T}$ 是同构. 命题~\bookref{7.4.62} 表明 $\underline{\Pic}_X$ 是 $\Bbbk$ 上局部有限型的代数空间, 且 (2) 成立. Picard 叠 $\underline{\mathcal P\!ic}_X$ 拟分离 (甚至有仿射对角), 故 $\underline{\Pic}_X$ 也拟分离. 这足以证明 $\underline{\Pic}_X$ 是分离概形, $\underline{\Pic}_X^0$ 拟射影. 这是定理~\bookref{5.5.28} 的直接结果, 但值得回顾其论证: $\underline{\Pic}_X$ 有稠密开子空间是概形 (定理~\bookref{5.5.1}), 群结构使我们可以平移这个开集, 从而证明 $\underline{\Pic}_X$ 是代数群, 它自动分离且连通分支拟射影 (命题~\bookref{B.1.19}). 这给出 (1).

对 (3), 只需证明 $\underline{\Pic}_X^0$ 固有. 已知它分离, 故只需验证固有性的赋值判据的存在性部分: 设 $R$ 是 $\Bbbk$ 上的 DVR, 分式域为 $K$, $L$ 是 $X_K$ 上的线丛. 由于 $X_R$ 正则, $L$ 延伸为 $X_R$ 上的线丛 $\widetilde L$ (例如, 若对除子 $D\subseteq X_K$ 有 $L=\mathcal O(D)$, 则取 $\widetilde L=\mathcal O(\overline D)$). (4) 中的光滑性来自特征 $0$ 的代数群光滑 (命题~\bookref{B.1.19}). 若 $L$ 是 $X$ 上的线丛, 则 Picard 叠与 Picard 概形的 Zariski 切空间一致, 形变理论 (命题~\bookref{C.1.18}) 表明 $T_{\underline{\Pic}_X,L}\cong\mathrm H^1(X,\mathcal O_X)$.
\end{proof}
\begin{remark}\label{src:7.4.66}
由 $\underline{\Pic}_X\cong\underline{\Pic}_X^{\mathrm{rel}}$ 的可表示性, $X\times\underline{\Pic}_X$ 上存在\emph{万有族} (或 \newterm{Poincaré 族}{Poincaré family}) $\mathcal P$, 满足: 对任意 $\Bbbk$-概形 $T$ 及 $X_T$ 上任意线丛 $L$, 存在唯一态射 $T\to\underline{\Pic}_X$, 使对某个 $T$ 上的线丛 $M$,
\[
L\cong\mathcal P|_{X\times T}\otimes p_2^*M.
\]
单位元的连通分支 $\underline{\Pic}_X^0$ 有如下函子描述: 它参数化与 $\mathcal O_X$ \emph{代数等价}的线丛 $L$ (即存在连通 $\Bbbk$-概形 $T$, 点 $t_0,t_1\in T(\Bbbk)$ 及 $X_T$ 上的线丛族 $\mathcal L$, 使 $\mathcal L_{t_0}\cong L$, $\mathcal L_{t_1}\cong\mathcal O_X$). 当 $X$ 是光滑曲线时, $\underline{\Pic}_X^0$ 参数化次数 $0$ 的线丛.
\end{remark}
\begin{remark}\label{src:7.4.67}
在特征 $p$ 下, Igusa 证明 $\underline{\Pic}(X)$ 可能不约化 \cite{Igu55}. 还应注意, 当 $X$ 不正规时 (例如结点或尖点曲线), $\underline{\Pic}^0(X)$ 不固有. Altman 与 Kleiman \cite{AK80} 通过分类秩 $1$ 的无挠层, 给出 $\underline{\Pic}^0(X)$ 的紧化. Picard 函子和概形有引人入胜的历史, 因为它们是代数几何中最早构造出的模空间之一. 关于这段历史的精彩介绍及 Picard 概形性质的更广泛讨论, 见 Kleiman 的文章 \cite{Kle05}.
\end{remark}
